\section*{Chapter \ref{ch:iv:online-assessments} --- Online Assessments}

\begin{solution}{iq:iv:online-assessments:1}
(a) $48 \times 0.35 = 48 \times \tfrac{7}{20} = 2.4 \times 7 = 16.8$. (b) $\tfrac78 - \tfrac23 = \tfrac{21 -
16}{24} = \tfrac{5}{24} \approx 0.2083$. (c) $2.4 \div 0.016 = 2400 \div 16 = 150$ (multiply both by 1000).
Each is one rewrite: a percentage as a fraction, a common denominator, a shift of the decimal point.
\iqlookfor{a rewrite that turns each item into one easy step.}
\end{solution}

\begin{solution}{iq:iv:online-assessments:2}
(a) 48: differences 5, 7, 9, 11 are odd numbers, so the next is 13 (the terms are $n^2 - 1$). (b) 33: each
term is twice the previous minus one. (c) 16: odd positions 1, 2, 3, 4 and even positions 4, 8, 12, so the next
even-position term is 16.
\iqlookfor{the method of differences, ratios and interleaving, tried in order, with a check.}
\end{solution}

\begin{solution}{iq:iv:online-assessments:3}
Answer when $p > w/(1 + w) = \tfrac13$. Blind: $p = \tfrac14$, worth $\tfrac14 - \tfrac12 \times \tfrac34 =
-\tfrac18$: skip. One option eliminated: $p = \tfrac13$, worth exactly 0: indifferent. Two eliminated: $p =
\tfrac12$, worth $+\tfrac14$: answer.
\iqlookfor{the break-even confidence of a scoring rule.}
\end{solution}

\begin{solution}{iq:iv:online-assessments:4}
One pass keeping the current run and the best: the run grows when a price exceeds the previous one and resets
to 1 otherwise; $O(n)$ time, $O(1)$ memory. Hidden cases: the empty list (answer 0), one price (1), all equal
(1, since strictly increasing), all decreasing (1), all increasing ($n$), equal neighbours inside a rise
(they break it), and $10^6$ prices (any $O(n^2)$ answer times out). The chapter's code checks the function
against a brute force on 300 random lists.
\iqlookfor{a linear scan, the edge cases named before coding, and the strictness of the inequality.}
\end{solution}

\begin{solution}{iq:iv:online-assessments:5}
The easy items take $60 \times 5 = 300$ seconds and are worth $0.97 - 0.03 = 0.94$ each, 56.4 in all. The
remaining 180 seconds allow 15 hard items, each worth $0.7 - 0.3 = 0.4$: 6 more, 62.4 in expectation. At 45\%
a hard item is worth $0.45 - 0.55 = -0.1$, so answer none of them and score 56.4; if an item looks easier than
its neighbours, answer it if you would give it better than even odds.
\iqlookfor{a time budget combined with the expected value of each answer.}
\end{solution}

\begin{solution}{iq:iv:online-assessments:6}
(a) The prime sums 2, 3, 5, 7, 11 occur 1, 2, 4, 6, 2 times out of 36: $15/36 = 5/12$. (b) By symmetry, at
least two heads out of three has probability $\tfrac12$ (it is the event that heads are the majority).
\iqlookfor{counting by the number of ways, and a symmetry argument where one exists.}
\end{solution}

\begin{solution}{iq:iv:online-assessments:7}
Seat 1 is C. D and E occupy two adjacent seats, D on the left. A is in seat 2, 3 or 4 and B is to A's right but
not next to A. If A is in seat 2, B is in seat 4 or 5; with D and E adjacent in the remaining seats this
forces A, D, E, B in seats 2 to 5 (B in 4 would leave seats 3 and 5 for D and E, which are not adjacent). If A
is in seat 3, B must be in seat 5 and D, E in seats 2 and 4, not adjacent: impossible. If A is in seat 4, B
has no seat. The row is C, A, D, E, B and D is in the middle; an exhaustive search over the 120 orders finds
this row alone.
\iqlookfor{placing the most constrained people first and eliminating cases on paper.}
\end{solution}

\begin{solution}{iq:iv:online-assessments:8}
With $k$ pumps the pot is $k$ if the burst point exceeds $k$, which happens with probability $(64 - k)/64$,
so the expected pot is $k(64 - k)/64$, maximised at $k = 32$ with value 16. Stopping at 20 gives 13.75. The
game probably also scores consistency (the same choice in the same situation), learning (adjusting when
bursts come earlier than expected) and the gap between the candidate's choice and the optimum, in both
directions: too little risk is as much a finding as too much.
\iqlookfor{an expected-value optimisation, and a view of what a behavioural game measures.}
\end{solution}

\begin{solution}{iq:iv:online-assessments:9}
Sort the prices once, $O(n \log n)$; answer each query with two binary searches, the first index not less
than $a$ and the first index greater than $b$, and return their difference: $O(\log n)$ per query, about
$4 \times 10^6$ steps in all, well inside the limit. A linear scan per query would be $4 \times 10^{10}$ and
time out. Edge cases: $a > b$ (answer 0), bounds equal to existing prices (the inclusive ends), duplicates, a
window outside all prices, and an empty list. The chapter's code checks it against a direct count on random
data.
\iqlookfor{preprocessing plus binary search, derived from the bounds, with inclusive-end care.}
\end{solution}

\begin{solution}{iq:iv:online-assessments:10}
One sitting: $\P(57 + 5Z > 60) = 1 - \Phi(0.6) \approx 0.274$. Two independent sittings: $1 - 0.726^2 \approx
0.473$. With practice, ability rises by $0.25 \times 8 = 2$ to 59: $1 - \Phi(0.2) \approx 0.421$ at one
sitting. Practice before the first attempt is worth almost as much as a second attempt, and a second attempt
is often not available for months (\cref{dat:iv:the-process-by-firm-type:published}): take the test when
practised, not when first invited, if the process allows a few days.
\iqlookfor{normal tail probabilities and the practical conclusion about timing.}
\end{solution}
